\section{Spiritual Knighthood}

\begin{quotex}
This is from a brief article by \textbf{Mircea Eliade} on the importance of \textbf{Ananda Coomaraswamy} and \textbf{Henry Corbin}. I don't know if this article is available elsewhere\footnote{\url{https://www.amiscorbin.com/images/documents/pdfs/Eliade_1979.pdf}}, but in case not, I have translated the French portions of it. 

\end{quotex}
Ananda Coomaraswamy did not try to prove any doctrine, but rather to show its consistency and intelligibility. Those properties provide the ground for faith, in the sense of ``confidence" rather than belief. AKC writes:

\begin{quotex}
[Philosophia Perennis] is neither a system nor a philosophy, it cannot be argued for or against. 

\end{quotex}
However, the story does not end there. Tradition cannot be proved because it is a fact, and there is no disputing facts. But to become a ``fact", it must be realized, i.e., made real. When this happens, it is no longer faith, but rather gnosis.

What Ananda Coomaraswamy did for the comparative study of the religions of India with those of the West, including ancient paganism and the spiritual system of the Medieval period, Henry Corbin did for what are referred to as the three ``Abrahamic faiths". He probes them in their depths, thereby rescuing them from the superficial critiques of both contemporary neo-atheism and neopaganism. Besides his lifelong devotion to the study of Persian and Sufi esoterism, Corbin was also interested in Heidegger's philosophy and, curiously, the immense works of Emmanuel Swedenborg.

Recognizing that the crisis of the West is in its foundations a spiritual crisis, Corbin founded the \textbf{International Center of Comparative Spiritual Research} in 1974. His goal was to encourage a new kind of university, where Traditional sciences could be studied free from the rationalistic prejudices of academia. He describes it:

\begin{quotex}
The specific character of this Association is marked in its very name: it makes the concept of ``spiritual knighthood" the norm of its researches and activities, and it makes Jerusalem the mystical ``symbol" of the encounters and gatherings whose place it hopes to be. The self-criticism of the West, as well as the accusations brought against it, never, in general, take into account the spiritual traditions of our western world, because they don't know them. The first cause of it is that unlike the great theological systems maintained by the religious Orders, or unlike the philosophical systems professed in the Universities, the treasure of the spiritual sciences, that can be grouped under the more or less suitable and adequate term of ``esoterism" is found abandoned in a state of neglect. One could just as well speak of suffocation by the canonical and legal spirit. The result is that this treasure remained buried in libraries, an object sometimes of the curiosity of well-intentioned scholars, but most often the prey of improvisers lacking discernment. Whence the proliferation of pseudo-esoterisms. It is therefore important to finally form a home for these high sciences whose neglect and oblivion are at once the cause and the symptom of the crisis of our civilization. To this purpose we cannot separate the history of philosophy, the history of the sciences, the history of spirituality. But it is possible to not separate them only by a ``rebirth" presupposing a plan of transhistoric permanency. Such is the meaning that we give to a restoration of traditional sciences and studies in the West. This restoration presupposed the necessary conjunction of the needs of the spiritual life and the rigors of scientific investigation, such that academics are accustomed to conduct. 

\end{quotex}
There are three stipulations: first, it is limited to the three ``religions of the book". Second, the spiritual meaning, i.e., the esoteric meaning of the Sacred Word is common to those religions. In his own words, this is the third:

\begin{quotex}
The Man, Adam, the \emph{Anthropos}, was created in some other parts, let us say in the Pleroma, and ``descended" down to this world. With him, the Word [Logos] descended to this world. With that descent, History begins. Man and the Word were made captive in or under a terrestrial envelope. Otherwise, they would not have been manifested in this world and History would never have begun. They would have remained in the state of flashing of non-perceptible light. However, under the terrestrial envelope, thanks to which we can see and hear, thanks to which we can give form to something like History, lives that flash of light which belongs to another world. At that point history is truly understood only if man perceives the trace of that flash and recognizes it … When he recognizes it, when that flash is remembered (\emph{anamnesis}), man experiences the state of ``those who know", the \emph{gnosis} of the gnostic, in the strict meaning of the word ``\emph{gnosis}"… we therefore say that the flash of light, \emph{exiled} under the terrestrial envelope, is at this point, \emph{saved}. And such is the deep meaning of the word \emph{gnosis}: a salvific knowledge because it is not a theoretical knowledge, but because it effects a transmutation of the inner man. It is the birth of the true man, the \emph{verus homo}. 

\end{quotex}
Corbin identified three sources of knowledge for ``traditional science":

\begin{enumerate}
\item \textbf{Philosophy}. Intellectual activity, in the metaphysical sense of \emph{nous} or \emph{intelectus}. 
\item \textbf{Theology}. The body of traditional teachings. 
\item \textbf{Theosophy}. Inner revelation, or active imagination. 
\end{enumerate}
For Ananda Coomaraswamy, the \emph{philosophia perennis} consists of the first principles behind the appearances of the several traditions. For Corbin, it is the visionary perception of the intermediary world that ``is designated in Arabic as ``\emph{alam al-mithal}", that I need to translate by \emph{mundus imaginalis}, the \emph{imaginal} world, to differentiate it from the \emph{imagination}.

Corbin concluded his presentation:

\begin{quotex}
Henceforth, the relationship between man and his God is that of knightly service. At its limit, this relationship produces the metamorphosis of the warrior knight to the spiritual knight.

I just tried to find the type of scholar who, in the measure where he corresponds to the idea of the legitimate heritage [the esoteric Abrahamic tradition], can assume the vocation of spiritual knight. Because he is at the same time the man of knowledge and the man of desire [yearning for the Absolute; see e.g., \emph{Man of Desire} by Louis Claude de St-Martin], because for him the intellectual life and research could never be isolated from the spiritual life and research, he provides a powerful contrast with the type of man in whom intellectuality is developed in ignorance of all spirituality. 

\end{quotex}


\flrightit{Posted on 2012-10-18 by Cologero }

\begin{center}* * *\end{center}

\begin{footnotesize}\begin{sffamily}



\texttt{Audrey on 2012-10-19 at 21:22 said: }

Indeed.. I had a brief moment of wondering `why' I do what I do.. and then it occurred to me that it largely is about boundaries. Without constant work, mechanicalness will win. Hence, I understood the notion of working with others.


\hfill

\texttt{Audrey on 2012-10-19 at 21:24 said: }

The mechanicalness of the Earth cannot be combatted with self-actulization only. Others must be brought along as well.


\end{sffamily}\end{footnotesize}
